% GENERATED FILE. Edit canonical lessons, then run npm run generate:books.
% canonical-source-hash: fnv1a64:d35616cccd217594
% canonical-lessons: GE-C360-vorurteil, GE-C360-toleranz, GE-C360-minderheit, GE-C360-diskriminieren, GE-C360-gleichberechtigung

\chapter{Vorurteile}
\label{ch:ge-vorurteile}
\hlchaptermodality{360}

\begin{chapteropening}
\textbf{By the end of this chapter:} \emph{I can talk about prejudice and tolerance, minorities, discrimination and equal rights.}

Complete the last lesson of chapter 360: I can talk about prejudice and tolerance, minorities, discrimination and equal rights.
\end{chapteropening}

\section[das Vorurteil]{das Vorurteil — the prejudice}

\label{lesson:GE-C360-vorurteil}



\begin{quote}
Before the new one: say the German for wealth, then the German for to put at a disadvantage.
\end{quote}

\subsection*{You'll want to know: das Vorurteil}
\textbf{das Vorurteil} — \textquotedblleft{}the prejudice\textquotedblright{}.

\begin{quote}
Er hat viele Vorurteile gegen Menschen aus der Stadt.
\end{quote}

\begin{quote}
\emph{He has a lot of prejudices against people from the city.}
\end{quote}

Plural: \emph{die Vorurteile}. Literally a judgement made beforehand.

\begin{cousinweb}
\emph{Vorurteil} joins \emph{vor} \textquotedblleft{}before\textquotedblright{} to \emph{Urteil} \textquotedblleft{}judgement\textquotedblright{}, from Old High German \emph{urteil} (EtymDB). It is usually explained as a rendering of Latin \emph{praeiūdicium} \textquotedblleft{}a judgement beforehand\textquotedblright{} (Lewis \& Short).
\end{cousinweb}

\begin{grammarlens}[title={What you've built}]
One more word for telling, explaining and describing at B1.
\end{grammarlens}

\subsection*{Your turn}
\begin{itemize}
\raggedright
  \item \emph{Say it:} \emph{das Vorurteil}
  \item \emph{Say it:} \emph{das Vorurteil}, once more
  \item \emph{Say it:} a sentence of your own with \emph{das Vorurteil}
  \item \emph{Recall:} say the German for wealth, then the German for to put at a disadvantage, then say \emph{das Vorurteil} again
\end{itemize}

\begin{tcolorbox}[breakable,colback=teal!4,colframe=teal!35!black,title={Before you move on}]
What does das Vorurteil mean? (\textquotedblleft{}The prejudice\textquotedblright{}.) Say it once more.
\end{tcolorbox}
\section[die Toleranz]{die Toleranz — tolerance}

\label{lesson:GE-C360-toleranz}



\begin{quote}
Before the new one: say the German for to put at a disadvantage, then the German for the prejudice.
\end{quote}

\subsection*{You'll want to know: die Toleranz}
\textbf{die Toleranz} — \textquotedblleft{}tolerance\textquotedblright{}.

\begin{quote}
In einer guten Schule lernen die Kinder Toleranz.
\end{quote}

\begin{quote}
\emph{In a good school the children learn tolerance.}
\end{quote}

No plural in this sense. The adjective is \emph{tolerant}.

\begin{cousinweb}
\emph{Toleranz} is a learned borrowing from Latin \emph{tolerantia} \textquotedblleft{}a bearing, endurance\textquotedblright{} (Lewis \& Short), from \emph{tolerāre} \textquotedblleft{}to bear, to put up with\textquotedblright{}.
\end{cousinweb}

\begin{grammarlens}[title={What you've built}]
One more word for telling, explaining and describing at B1.
\end{grammarlens}

\subsection*{Your turn}
\begin{itemize}
\raggedright
  \item \emph{Say it:} \emph{die Toleranz}
  \item \emph{Say it:} \emph{die Toleranz}, once more
  \item \emph{Say it:} a sentence of your own with \emph{die Toleranz}
  \item \emph{Recall:} say the German for to put at a disadvantage, then the German for the prejudice, then say \emph{die Toleranz} again
\end{itemize}

\begin{tcolorbox}[breakable,colback=teal!4,colframe=teal!35!black,title={Before you move on}]
What does die Toleranz mean? (\textquotedblleft{}Tolerance\textquotedblright{}.) Say it once more.
\end{tcolorbox}
\section[die Minderheit]{die Minderheit — the minority}

\label{lesson:GE-C360-minderheit}



\begin{quote}
Before the new one: say the German for the prejudice, then the German for tolerance.
\end{quote}

\subsection*{You'll want to know: die Minderheit}
\textbf{die Minderheit} — \textquotedblleft{}the minority\textquotedblright{}.

\begin{quote}
In diesem Land gibt es mehrere sprachliche Minderheiten.
\end{quote}

\begin{quote}
\emph{There are several language minorities in this country.}
\end{quote}

Plural: \emph{die Minderheiten}. The opposite is \emph{die Mehrheit}.

\begin{cousinweb}
\emph{Minderheit} is built from \emph{minder} \textquotedblleft{}less\textquotedblright{}, inherited from Middle High German \emph{minder}, Old High German \emph{minniro} and Proto-Germanic \emph{minnizô} (EtymDB), a cousin of Latin \emph{minor}, with \emph{-heit}.
\end{cousinweb}

\begin{grammarlens}[title={What you've built}]
One more word for telling, explaining and describing at B1.
\end{grammarlens}

\subsection*{Your turn}
\begin{itemize}
\raggedright
  \item \emph{Say it:} \emph{die Minderheit}
  \item \emph{Say it:} \emph{die Minderheit}, once more
  \item \emph{Say it:} a sentence of your own with \emph{die Minderheit}
  \item \emph{Recall:} say the German for the prejudice, then the German for tolerance, then say \emph{die Minderheit} again
\end{itemize}

\begin{tcolorbox}[breakable,colback=teal!4,colframe=teal!35!black,title={Before you move on}]
What does die Minderheit mean? (\textquotedblleft{}The minority\textquotedblright{}.) Say it once more.
\end{tcolorbox}
\section[diskriminieren]{diskriminieren — to discriminate against}

\label{lesson:GE-C360-diskriminieren}



\begin{quote}
Before the new one: say the German for tolerance, then the German for the minority.
\end{quote}

\subsection*{You'll want to know: diskriminieren}
\textbf{diskriminieren} — \textquotedblleft{}to discriminate against\textquotedblright{}.

\begin{quote}
Die Firma darf niemanden wegen seiner Herkunft diskriminieren.
\end{quote}

\begin{quote}
\emph{The firm may not discriminate against anyone because of their origin.}
\end{quote}

Perfekt: \emph{hat diskriminiert}. It takes a direct object: \emph{jemanden diskriminieren}.

\begin{cousinweb}
\emph{diskriminieren} is a learned borrowing from Latin \emph{discrīmināre} \textquotedblleft{}to divide, to separate\textquotedblright{}, from \emph{discrīmen} \textquotedblleft{}that which divides\textquotedblright{} (Lewis \& Short): you set people apart.
\end{cousinweb}

\begin{grammarlens}[title={What you've built}]
One more word for telling, explaining and describing at B1.
\end{grammarlens}

\subsection*{Your turn}
\begin{itemize}
\raggedright
  \item \emph{Say it:} \emph{diskriminieren}
  \item \emph{Say it:} \emph{diskriminieren}, once more
  \item \emph{Say it:} a sentence of your own with \emph{diskriminieren}
  \item \emph{Recall:} say the German for tolerance, then the German for the minority, then say \emph{diskriminieren} again
\end{itemize}

\begin{tcolorbox}[breakable,colback=teal!4,colframe=teal!35!black,title={Before you move on}]
What does diskriminieren mean? (\textquotedblleft{}To discriminate against\textquotedblright{}.) Say it once more.
\end{tcolorbox}
\section[die Gleichberechtigung]{die Gleichberechtigung — equal rights, equality}

\label{lesson:GE-C360-gleichberechtigung}



\begin{quote}
Before the new one: say the German for the minority, then the German for to discriminate against.
\end{quote}

\subsection*{You'll want to know: die Gleichberechtigung}
\textbf{die Gleichberechtigung} — \textquotedblleft{}equal rights, equality\textquotedblright{}.

\begin{quote}
Die Gleichberechtigung von Männern und Frauen ist noch nicht überall Realität.
\end{quote}

\begin{quote}
\emph{Equal rights for men and women are not yet a reality everywhere.}
\end{quote}

No plural. The adjective is \emph{gleichberechtigt} — with equal rights.

\begin{cousinweb}
\emph{Gleichberechtigung} joins \emph{gleich} \textquotedblleft{}equal\textquotedblright{}, Proto-Germanic \emph{galīkaz} (EtymDB), to \emph{Berechtigung} \textquotedblleft{}entitlement\textquotedblright{}, from \emph{Recht} \textquotedblleft{}right\textquotedblright{}, Proto-Germanic \emph{rehtaz} (EtymDB).
\end{cousinweb}

\begin{grammarlens}[title={What you've built}]
One more word for telling, explaining and describing at B1.
\end{grammarlens}

\subsection*{Your turn}
\begin{itemize}
\raggedright
  \item \emph{Say it:} \emph{die Gleichberechtigung}
  \item \emph{Say it:} \emph{die Gleichberechtigung}, once more
  \item \emph{Say it:} a sentence of your own with \emph{die Gleichberechtigung}
  \item \emph{Recall:} say the German for the minority, then the German for to discriminate against, then say \emph{die Gleichberechtigung} again
\end{itemize}

\begin{tcolorbox}[breakable,colback=teal!4,colframe=teal!35!black,title={Before you move on}]
What does die Gleichberechtigung mean? (\textquotedblleft{}Equal rights, equality\textquotedblright{}.) Say it once more.
\end{tcolorbox}
